\citet{uesato2022} compare feedback on intermediate steps with feedback on final answers in mathematical reasoning. \citet{lightman2023} study process supervision on MATH and release step-level annotations. These studies motivate evaluating intermediate reasoning rather than relying only on the final answer.

Recent benchmarks examine how well evaluators detect reasoning errors. ProcessBench focuses on identifying the earliest incorrect step in a solution \citep{zheng2025}. PRMBench evaluates process-level reward models through dimensions including simplicity, soundness, and sensitivity \citep{song2025}. Its Confidence Invariance category specifically assesses whether PRMs can detect incorrect reasoning steps expressed with high confidence.

Research on human preferences raises a related concern. \citet{sharma2023} find that human evaluators and preference models sometimes prefer persuasive responses that agree with a user's mistaken belief over truthful responses. This raises the question of whether evaluation scores depend on how a response is written as well as whether it is correct. Their study examines preferences for whole responses. It does not test PRM scores or the effect of adding a fixed confidence phrase to a mathematical step. Our study focuses on within-question changes in PRM scores and correct–incorrect score gaps when a fixed confident phrase replaces a neutral one.
